% ============================================================
% C2 Challenge: AI for Math  --  paper.tex
% From Natural Language to Verifiable Reasoning:
% A Semantic Reduction Framework for Reliable AI Mathematics
%
% 编译：pdflatex paper.tex && bibtex paper && pdflatex paper.tex && pdflatex paper.tex
% 要求：可独立编译，引用可解析。references.bib 中 11 篇全部经人工核验。
% ============================================================

\documentclass[11pt,a4paper]{article}

\usepackage[utf8]{inputenc}
\usepackage[T1]{fontenc}
\usepackage[margin=1in]{geometry}
\usepackage{amsmath,amssymb}
\usepackage{graphicx}
\usepackage{booktabs}
\usepackage{hyperref}

\title{From Natural Language to Verifiable Reasoning:\\
A Semantic Reduction Framework for Reliable AI Mathematics}

\author{Feng Shuo}

\date{\today}

\begin{document}
\maketitle

% ------------------------------------------------------------
\begin{abstract}
Large language models produce mathematical arguments that are fluent but not
trustworthy: a proof may read correctly and still be wrong.  We argue that this
failure is not primarily a failure of generation, nor of verification, but of
the \emph{middle layer} that connects them.  We decompose an AI4Math system into
three layers --- generation, semantic reduction, and verification --- and show
that the second layer is the actual bottleneck: verification can only operate on
computable objects, and natural language is not one.  We propose a structured
typed intermediate representation (Typed IR) as the core system component, and
sketch an evaluation protocol that reports \emph{verifiability} alongside
accuracy.  We do not claim a new solver; we claim a sharper decomposition and a
measurable target.

\noindent\textbf{Keywords:} AI for mathematics, formal verification, semantic
reduction, neuro-symbolic systems, reliability
\end{abstract}

% ------------------------------------------------------------
\section{Introduction}

Machine learning has entered core mathematical practice.  Neural models guided
human intuition toward new results in knot theory and representation
theory~\cite{davies2021advancing}; an evolutionary procedure paired with a
program evaluator discovered new cap-set constructions beyond the best known
ones~\cite{romera2024funsearch}; and a neuro-symbolic system solved 25 of 30
olympiad geometry problems~\cite{trinh2024alphageometry}.

Yet these successes sit next to an uncomfortable observation.  Language models
regularly emit reasoning chains that are locally plausible and globally
wrong~\cite{wei2022chain}.  On competition-level benchmarks such as
MATH~\cite{hendrycks2021math} and GSM8K~\cite{cobbe2021gsm8k}, a correct final
answer does not certify a correct derivation.  Process supervision, which labels
intermediate steps rather than only outcomes, improves
reliability~\cite{lightman2023lets} --- but it is expensive precisely because
there is no machine-checkable handle on an intermediate step written in natural
language.

\textbf{Claim.}  Reliability cannot be bought by final-answer verification
alone.  It requires an explicit \emph{semantic reduction layer} that maps
natural-language reasoning into a structured object a verifier can act on.

\paragraph{Contributions.}
\begin{enumerate}
  \item A three-layer decomposition of AI4Math systems that isolates semantic
        reduction as a first-class component.
  \item An argument that most observed failures concentrate in that layer, not
        in generation or verification.
  \item A proposed Typed IR design with bidirectional consistency requirements,
        plus an evaluation protocol reporting verifiability in addition to
        accuracy.
\end{enumerate}

% ------------------------------------------------------------
\section{A Three-Layer Decomposition}

Let an AI4Math system be decomposed as $G \to R \to V$:

\subsection{Layer 1: Generation}
A language model emits a natural-language reasoning chain, typically via
chain-of-thought prompting~\cite{wei2022chain}.  Its output is
probabilistic, untyped, and semantically underspecified.

\subsection{Layer 2: Semantic Reduction}
Natural language is mapped to a structured representation: first- or
higher-order logic expressions, executable programs, proof objects, or a typed
intermediate representation.  The theoretical anchor is the Curry--Howard
correspondence.  Three failure modes dominate here:
\begin{itemize}
  \item \textbf{Ambiguity}: the same sentence admits several formal readings.
  \item \textbf{Missing typing}: implicit domains, coercions, and side
        conditions are silently dropped.
  \item \textbf{Non-invertibility}: the mapping loses information, so a
        verified object cannot be faithfully rendered back.
\end{itemize}

\subsection{Layer 3: Verification}
A formal system checks self-consistency, derivability, factual grounding, and
constraint satisfaction.  In practice this is an interactive theorem prover such
as Lean~4~\cite{moura2021lean4} over the mathlib
library~\cite{mathlib2020}, or a symbolic/SMT solver.  The cross-system
benchmark miniF2F~\cite{zheng2022minif2f} shows that this layer is now
measurable and comparable across provers.

% ------------------------------------------------------------
\section{The Semantic Reduction Bottleneck}

We claim that when an AI4Math system fails, it fails in Layer~2 more often than
in Layers~1 or~3.

The argument is structural.  Verification is only defined over objects with a
decidable notion of well-formedness.  Natural language has none.  Therefore, if
Layer~2 produces no object, or produces an object that does not faithfully encode
the intended claim, then Layer~3 has nothing to check --- or checks the wrong
thing.  Two consequences follow:
\begin{enumerate}
  \item \textbf{Verification failure}: the pipeline reports ``unverified''
        rather than ``wrong'', conflating a translation failure with a
        reasoning failure.
  \item \textbf{Undetectable hallucination}: an unfaithful translation can be
        internally consistent and still prove a proposition different from the
        one that was asked.
\end{enumerate}

This reframes the reliability problem: the scarce resource is not the prover's
power but the \emph{fidelity of the translation}.

% ------------------------------------------------------------
\section{Proposed Framework}

\subsection{Structured Intermediate Representation}
We propose a typed IR with four components: a type system, an operational
semantics, a constraint system, and a composable graph/AST structure.

\subsection{Pipeline}
The standard flow becomes
\[
\text{Natural language} \;\to\; \text{Typed IR} \;\to\;
\text{Formal verification} \;\to\; \text{Natural-language rendering},
\]
where natural language is the \emph{interface} and the typed IR is the
\emph{kernel}.

\subsection{Bidirectional Consistency}
We require $\text{IR} \leftrightarrow \text{natural language}$ to be both
interpretable and traceable: every IR node records the span it came from, and
every rendered sentence records the node it came from.  This is the property
that makes a human audit possible.

% ------------------------------------------------------------
\section{Related Work}

\textbf{Generation-side.}  Chain-of-thought~\cite{wei2022chain} and
process supervision~\cite{lightman2023lets} improve step quality but leave the
steps unverifiable.

\textbf{Verification-side.}  Neuro-symbolic theorem
proving~\cite{trinh2024alphageometry,zheng2022minif2f} and the Lean/mathlib
ecosystem~\cite{moura2021lean4,mathlib2020} provide sound kernels, but assume a
formal statement already exists.

\textbf{Search-side.}  Program-search methods~\cite{romera2024funsearch} and
ML-guided intuition~\cite{davies2021advancing} sidestep verification by using
an evaluator as a filter --- effective where an evaluator exists, silent where
it does not.

\textbf{Difference.}  We treat semantic reduction as the core system component,
not as preprocessing, and we make its fidelity a reported metric.

% ------------------------------------------------------------
\section{Experimental Design}

\textbf{Benchmarks.}  GSM8K~\cite{cobbe2021gsm8k} and
MATH~\cite{hendrycks2021math} for informal reasoning; miniF2F~\cite{zheng2022minif2f}
for formal proving.

\textbf{Comparison.}  Pure LLM baseline versus IR-based pipeline.

\textbf{Metrics.}  Beyond accuracy we report:
\begin{itemize}
  \item \emph{Verifiability rate}: fraction of outputs that yield a
        machine-checkable object.
  \item \emph{Consistency}: agreement between the rendered statement and the
        verified statement.
\end{itemize}

We note honestly that this design is a proposal, not a completed experiment;
running it requires engineering effort beyond the scope of the present
write-up.

% ------------------------------------------------------------
\section{Discussion}

\textbf{Limitations.}  IR design is expensive; semantic mapping cost may exceed
the cost of the proof it enables; and the framework is silent on conjectures
that cannot be stated in the target logic at all.

\textbf{Extensions.}  The same decomposition should apply to general reasoning
systems and to agent architectures, where the ``verifier'' may be an external
tool rather than a proof kernel.

% ------------------------------------------------------------
\section{Conclusion}

Reliable AI reasoning requires generation, semantic reduction, and
verification together.  The contribution is to name the middle term, argue that
it is the bottleneck, and propose a measurable notion of verifiability that
makes the claim testable rather than rhetorical.

% ------------------------------------------------------------
\bibliographystyle{plain}
\bibliography{references}

\end{document}
